Memory allocation is very efficient in garbage collectors used in
state-of-the-art compilers such as ocamlopt and ghc. The OCaml garbage
collector is a modern hybrid generational and incremental collector
that is known to often outperform manual memory
allocation. Generational garbage collectors also improve cache
locality~\cite{ZornGC} which is an important feature for modern
hardware architectures where poor locality often reduces program
performance~\cite{cachelocality}.

We initially used the standard \emph{malloc} and \emph{free} from the C standard library to implement
our runtime, and have observed poor performance in many benchmarks.
Profiling data made it clear that \emph{malloc} and \emph{free} were the bottlenecks.
We addressed this issue by switching to \emph{jemalloc}~\cite{jemalloc}, a \emph{malloc}
and \emph{free} replacement, and have observed a significant performance
improvement in many benchmarks by up to 3x. We obtained another $30\%$
performance improvement by implementing our own custom memory allocator
optimized for small values. Our allocator uses a new memory allocation
technique developed by \citet{freelistsharding} for the Koka programming
language~\cite{Koka} runtime.

We have also implemented lazy reference counting as a configuration option
in our runtime for users that might be concerned with long pauses when
deallocating values. The idea is quite simple: The instruction
$\kw{dec}\ x$ just puts $x$ in a deletion queue if its reference count
becomes $0$.  Before we allocate memory, we check the deletion queue,
and remove one value if it is not empty, and invoke $\kw{dec}$ on
nested values. \citet{BoehmCostRC} shows that lazy reference
counting may result in an unacceptable worst-case space overhead. We
did not manage to observe the worst-case behavior in any of our
benchmarks, but further investigation is needed. We suspect that
providing a primitive for (partially) flushing the lazy reference
counting queue should be sufficient to address any program that may
experience a significant space overhead.
